% ============================================================================
% 引言 - 01_introduction.tex
% ============================================================================

\section{引言}
\label{sec:introduction}
（研究背景）城市低空无人机应用增长，但公共安全和社会接受度成为瓶颈。
（现有研究进展）现有路径规划多关注距离、能耗、静态障碍或静态风险，难以描述城市低空风险的动态性。
（核心研究缺口）指出风险有两层：失效概率是动态的，事故后果也是动态的；同时噪声是正常运行的确定性社会成本。
（本文研究问题）构建时空风险张量，并用 EDA-Kmeans-CostA* 高效求解。
（本文核心贡献）列出贡献。因此，本文首次系统性地建立了一个时空耦合的动态 TPR 模型（把你的气象动态因子、人口潮汐引力、噪声时间惩罚全部抛出来）。

城市~\cite{ref1, ref2}。



本文其余部分组织如下：第~\ref{sec:related_work}节回顾相关工作并指出现有研究局限。
第~\ref{sec:tpr_models}节介绍第三方风险成本评估模型。第~\ref{sec:problem_formulation}节形式化综合风险成本函数并定义优化问题。
第~\ref{sec:algorithm}节描述我们的时空风险感知路径规划算法。
第~\ref{sec:case_study}节展示城市场景案例研究和实验分析。
最后，第~\ref{sec:conclusion}节总结全文并概述未来研究方向。